\chapter{Introduction}

\section{Background and Overview}
The global financial landscape has experienced an unprecedented transition from physical, branch-centric operations to ubiquitous digital banking platforms. Modern users demand instant, frictionless financial interactions accessible across various computing devices. However, this convenience introduces substantial engineering hurdles: ensuring ironclad security for sensitive customer assets, preventing unauthorized transaction replay, maintaining continuous availability, and delivering intuitive, responsive user interfaces.

Traditional financial platforms often struggle with legacy monolithic designs, brittle relational database dependencies, and cluttered user interfaces that introduce user fatigue. \textbf{SecureBank} addresses these challenges by reimagining digital banking through a decoupled client-server architecture built on modern JavaScript/TypeScript paradigms. The system combines an ultra-modern aesthetic inspired by the Aceternity Studio design philosophy with a hardened Node.js/Express application core backed by a resilient dual-engine persistence layer.

\section{Motivation}
The impetus behind developing SecureBank stems from three primary observations in contemporary financial software engineering:
\begin{enumerate}
    \item \textbf{Vulnerability of Standard Web Systems:} Many simulation or academic banking frameworks store passwords in plaintext or employ obsolete hashing algorithms (such as MD5 or SHA-1) vulnerable to rainbow table attacks and collision exploitation. There is a pressing need for reference implementations utilizing salted adaptive key derivation (Bcrypt) and secure HTTP cookie architectures.
    \item \textbf{Brittle Database Coupling:} Many applications crash upon initialization if their primary cloud or distributed database cluster experiences network partitions or offline developer scenarios. An engineering pattern that provides transparent fallback to a fast embedded engine (SQLite) preserves developer productivity while ensuring cloud scalability (MongoDB).
    \item \textbf{Cognitive Friction in Financial Interfaces:} Traditional banking interfaces suffer from visual bloat and disorganized information architecture. Designing a high-trust interface using architectural framing, clear typography, and responsive state management establishes immediate user confidence.
\end{enumerate}

\section{Problem Statement}
To design, implement, and deploy an enterprise-ready digital banking simulation system that provides:
\begin{itemize}
    \item Robust customer onboarding with strict cryptographic password hashing and schema validation.
    \item Synchronous double-entry transaction processing with atomic balance enforcement to eliminate drift.
    \item An administrative governance portal for real-time liquidity management, customer monitoring, and global audit logging.
    \item A resilient data persistence pipeline supporting MongoDB as the primary document store with automatic, crash-free fallback to an embedded SQLite database.
    \item A cloud-ready deployment configuration suitable for modern serverless container execution (Vercel).
\end{itemize}

\section{Project Objectives}
The core objectives of the SecureBank project are as follows:
\begin{enumerate}
    \item \textbf{Cryptographic Security:} Implement 12-round salted Bcrypt password encryption, secure session regeneration on privilege elevation, and \texttt{httpOnly} cookie transmission to mitigate XSS and CSRF threats.
    \item \textbf{Dual-Engine Storage Abstraction:} Construct a unified Data Access Layer (\texttt{dbService}) that automatically routes operations to MongoDB (via Mongoose) when connected, or falls back to an embedded SQLite database without throwing unhandled connection exceptions.
    \item \textbf{Atomic Transaction Enforcement:} Guarantee that balance recalculations and transaction ledger entries occur as atomic units, ensuring zero ledger drift during simulated debits or credits.
    \item \textbf{Administrative Control Center:} Provide an authorized administrative tier capable of querying system metrics, viewing all registered accounts, adjusting liquidity, and purging user accounts with cascading ledger cleanup.
    \item \textbf{Modern High-Trust Frontend:} Implement a responsive React 19 SPA incorporating floating pill navigation, bento-grid layouts, dynamic virtual DOM updates, and custom CSS styling without external bloated UI component kits.
    \item \textbf{Continuous Deployment:} Configure the platform for automated cloud hosting via Vercel serverless lambda wrappers and GitHub synchronization.
\end{enumerate}

\section{Scope of the Project}
The scope of this project encompasses:
\begin{itemize}
    \item \textbf{User Module:} Customer self-registration, validation of 10-digit Indian mobile formats, ISO birthdates, unique email verification, secure session-based authentication, and account profile deletion.
    \item \textbf{Transaction Module:} Simulated debit and credit operations with instantaneous validation against available balances, generating immutable chronological activity records.
    \item \textbf{Administration Module:} Dedicated admin portal protected by role-based guards, supporting customer search, balance intervention, audit inspection, and account termination.
    \item \textbf{Database Module:} Mongoose schema definitions for MongoDB and table structures for SQLite with automatic data migration from local storage to cloud MongoDB upon initial synchronization.
\end{itemize}
\textit{Note on Exclusions:} SecureBank is an engineering prototype and simulation environment. It does not interface with actual Central Bank payment gateways (e.g., UPI, NEFT, RTGS) or real monetary assets.

\section{Organization of the Report}
The subsequent chapters of this report are structured as follows:
\begin{itemize}
    \item \textbf{Chapter 2 (Literature Survey):} Examines related work, existing open-source and commercial banking architectures, and highlights identified research and technical gaps.
    \item \textbf{Chapter 3 (System Requirements and Analysis):} Outlines functional and non-functional requirements, user personas, software/hardware prerequisites, and feasibility dimensions.
    \item \textbf{Chapter 4 (System Design and Architecture):} Details the overall multi-tier architectural framework, Data Flow Diagrams (DFDs), UML diagrams, and entity-relationship models.
    \item \textbf{Chapter 5 (Implementation Details):} Discusses the core algorithmic implementation, cryptographic routines, dual-engine data service logic, API endpoints, and frontend components.
    \item \textbf{Chapter 6 (Results, Testing, and Discussion):} Presents verification results, testing methodologies (functional, validation, security), comparative analysis, and concludes with future enhancements.
    \item \textbf{Appendix \& Bibliography:} Provides the API route specification, configuration instructions, and academic references.
\end{itemize}
